% Options for packages loaded elsewhere
% Options for packages loaded elsewhere
\PassOptionsToPackage{unicode}{hyperref}
\PassOptionsToPackage{hyphens}{url}
\PassOptionsToPackage{dvipsnames,svgnames,x11names}{xcolor}
%
\documentclass[
  letterpaper,
  DIV=11,
  numbers=noendperiod]{scrartcl}
\usepackage{xcolor}
\usepackage{amsmath,amssymb}
\setcounter{secnumdepth}{-\maxdimen} % remove section numbering
\usepackage{iftex}
\ifPDFTeX
  \usepackage[T1]{fontenc}
  \usepackage[utf8]{inputenc}
  \usepackage{textcomp} % provide euro and other symbols
\else % if luatex or xetex
  \usepackage{unicode-math} % this also loads fontspec
  \defaultfontfeatures{Scale=MatchLowercase}
  \defaultfontfeatures[\rmfamily]{Ligatures=TeX,Scale=1}
\fi
\usepackage{lmodern}
\ifPDFTeX\else
  % xetex/luatex font selection
\fi
% Use upquote if available, for straight quotes in verbatim environments
\IfFileExists{upquote.sty}{\usepackage{upquote}}{}
\IfFileExists{microtype.sty}{% use microtype if available
  \usepackage[]{microtype}
  \UseMicrotypeSet[protrusion]{basicmath} % disable protrusion for tt fonts
}{}
\makeatletter
\@ifundefined{KOMAClassName}{% if non-KOMA class
  \IfFileExists{parskip.sty}{%
    \usepackage{parskip}
  }{% else
    \setlength{\parindent}{0pt}
    \setlength{\parskip}{6pt plus 2pt minus 1pt}}
}{% if KOMA class
  \KOMAoptions{parskip=half}}
\makeatother
% Make \paragraph and \subparagraph free-standing
\makeatletter
\ifx\paragraph\undefined\else
  \let\oldparagraph\paragraph
  \renewcommand{\paragraph}{
    \@ifstar
      \xxxParagraphStar
      \xxxParagraphNoStar
  }
  \newcommand{\xxxParagraphStar}[1]{\oldparagraph*{#1}\mbox{}}
  \newcommand{\xxxParagraphNoStar}[1]{\oldparagraph{#1}\mbox{}}
\fi
\ifx\subparagraph\undefined\else
  \let\oldsubparagraph\subparagraph
  \renewcommand{\subparagraph}{
    \@ifstar
      \xxxSubParagraphStar
      \xxxSubParagraphNoStar
  }
  \newcommand{\xxxSubParagraphStar}[1]{\oldsubparagraph*{#1}\mbox{}}
  \newcommand{\xxxSubParagraphNoStar}[1]{\oldsubparagraph{#1}\mbox{}}
\fi
\makeatother

\usepackage{color}
\usepackage{fancyvrb}
\newcommand{\VerbBar}{|}
\newcommand{\VERB}{\Verb[commandchars=\\\{\}]}
\DefineVerbatimEnvironment{Highlighting}{Verbatim}{commandchars=\\\{\}}
% Add ',fontsize=\small' for more characters per line
\usepackage{framed}
\definecolor{shadecolor}{RGB}{241,243,245}
\newenvironment{Shaded}{\begin{snugshade}}{\end{snugshade}}
\newcommand{\AlertTok}[1]{\textcolor[rgb]{0.68,0.00,0.00}{#1}}
\newcommand{\AnnotationTok}[1]{\textcolor[rgb]{0.37,0.37,0.37}{#1}}
\newcommand{\AttributeTok}[1]{\textcolor[rgb]{0.40,0.45,0.13}{#1}}
\newcommand{\BaseNTok}[1]{\textcolor[rgb]{0.68,0.00,0.00}{#1}}
\newcommand{\BuiltInTok}[1]{\textcolor[rgb]{0.00,0.23,0.31}{#1}}
\newcommand{\CharTok}[1]{\textcolor[rgb]{0.13,0.47,0.30}{#1}}
\newcommand{\CommentTok}[1]{\textcolor[rgb]{0.37,0.37,0.37}{#1}}
\newcommand{\CommentVarTok}[1]{\textcolor[rgb]{0.37,0.37,0.37}{\textit{#1}}}
\newcommand{\ConstantTok}[1]{\textcolor[rgb]{0.56,0.35,0.01}{#1}}
\newcommand{\ControlFlowTok}[1]{\textcolor[rgb]{0.00,0.23,0.31}{\textbf{#1}}}
\newcommand{\DataTypeTok}[1]{\textcolor[rgb]{0.68,0.00,0.00}{#1}}
\newcommand{\DecValTok}[1]{\textcolor[rgb]{0.68,0.00,0.00}{#1}}
\newcommand{\DocumentationTok}[1]{\textcolor[rgb]{0.37,0.37,0.37}{\textit{#1}}}
\newcommand{\ErrorTok}[1]{\textcolor[rgb]{0.68,0.00,0.00}{#1}}
\newcommand{\ExtensionTok}[1]{\textcolor[rgb]{0.00,0.23,0.31}{#1}}
\newcommand{\FloatTok}[1]{\textcolor[rgb]{0.68,0.00,0.00}{#1}}
\newcommand{\FunctionTok}[1]{\textcolor[rgb]{0.28,0.35,0.67}{#1}}
\newcommand{\ImportTok}[1]{\textcolor[rgb]{0.00,0.46,0.62}{#1}}
\newcommand{\InformationTok}[1]{\textcolor[rgb]{0.37,0.37,0.37}{#1}}
\newcommand{\KeywordTok}[1]{\textcolor[rgb]{0.00,0.23,0.31}{\textbf{#1}}}
\newcommand{\NormalTok}[1]{\textcolor[rgb]{0.00,0.23,0.31}{#1}}
\newcommand{\OperatorTok}[1]{\textcolor[rgb]{0.37,0.37,0.37}{#1}}
\newcommand{\OtherTok}[1]{\textcolor[rgb]{0.00,0.23,0.31}{#1}}
\newcommand{\PreprocessorTok}[1]{\textcolor[rgb]{0.68,0.00,0.00}{#1}}
\newcommand{\RegionMarkerTok}[1]{\textcolor[rgb]{0.00,0.23,0.31}{#1}}
\newcommand{\SpecialCharTok}[1]{\textcolor[rgb]{0.37,0.37,0.37}{#1}}
\newcommand{\SpecialStringTok}[1]{\textcolor[rgb]{0.13,0.47,0.30}{#1}}
\newcommand{\StringTok}[1]{\textcolor[rgb]{0.13,0.47,0.30}{#1}}
\newcommand{\VariableTok}[1]{\textcolor[rgb]{0.07,0.07,0.07}{#1}}
\newcommand{\VerbatimStringTok}[1]{\textcolor[rgb]{0.13,0.47,0.30}{#1}}
\newcommand{\WarningTok}[1]{\textcolor[rgb]{0.37,0.37,0.37}{\textit{#1}}}

\usepackage{longtable,booktabs,array}
\usepackage{calc} % for calculating minipage widths
% Correct order of tables after \paragraph or \subparagraph
\usepackage{etoolbox}
\makeatletter
\patchcmd\longtable{\par}{\if@noskipsec\mbox{}\fi\par}{}{}
\makeatother
% Allow footnotes in longtable head/foot
\IfFileExists{footnotehyper.sty}{\usepackage{footnotehyper}}{\usepackage{footnote}}
\makesavenoteenv{longtable}
\usepackage{graphicx}
\makeatletter
\newsavebox\pandoc@box
\newcommand*\pandocbounded[1]{% scales image to fit in text height/width
  \sbox\pandoc@box{#1}%
  \Gscale@div\@tempa{\textheight}{\dimexpr\ht\pandoc@box+\dp\pandoc@box\relax}%
  \Gscale@div\@tempb{\linewidth}{\wd\pandoc@box}%
  \ifdim\@tempb\p@<\@tempa\p@\let\@tempa\@tempb\fi% select the smaller of both
  \ifdim\@tempa\p@<\p@\scalebox{\@tempa}{\usebox\pandoc@box}%
  \else\usebox{\pandoc@box}%
  \fi%
}
% Set default figure placement to htbp
\def\fps@figure{htbp}
\makeatother


% definitions for citeproc citations
\NewDocumentCommand\citeproctext{}{}
\NewDocumentCommand\citeproc{mm}{%
  \begingroup\def\citeproctext{#2}\cite{#1}\endgroup}
\makeatletter
 % allow citations to break across lines
 \let\@cite@ofmt\@firstofone
 % avoid brackets around text for \cite:
 \def\@biblabel#1{}
 \def\@cite#1#2{{#1\if@tempswa , #2\fi}}
\makeatother
\newlength{\cslhangindent}
\setlength{\cslhangindent}{1.5em}
\newlength{\csllabelwidth}
\setlength{\csllabelwidth}{3em}
\newenvironment{CSLReferences}[2] % #1 hanging-indent, #2 entry-spacing
 {\begin{list}{}{%
  \setlength{\itemindent}{0pt}
  \setlength{\leftmargin}{0pt}
  \setlength{\parsep}{0pt}
  % turn on hanging indent if param 1 is 1
  \ifodd #1
   \setlength{\leftmargin}{\cslhangindent}
   \setlength{\itemindent}{-1\cslhangindent}
  \fi
  % set entry spacing
  \setlength{\itemsep}{#2\baselineskip}}}
 {\end{list}}
\usepackage{calc}
\newcommand{\CSLBlock}[1]{\hfill\break\parbox[t]{\linewidth}{\strut\ignorespaces#1\strut}}
\newcommand{\CSLLeftMargin}[1]{\parbox[t]{\csllabelwidth}{\strut#1\strut}}
\newcommand{\CSLRightInline}[1]{\parbox[t]{\linewidth - \csllabelwidth}{\strut#1\strut}}
\newcommand{\CSLIndent}[1]{\hspace{\cslhangindent}#1}



\setlength{\emergencystretch}{3em} % prevent overfull lines

\providecommand{\tightlist}{%
  \setlength{\itemsep}{0pt}\setlength{\parskip}{0pt}}



 


\KOMAoption{captions}{tableheading}
\makeatletter
\@ifpackageloaded{tcolorbox}{}{\usepackage[skins,breakable]{tcolorbox}}
\@ifpackageloaded{fontawesome5}{}{\usepackage{fontawesome5}}
\definecolor{quarto-callout-color}{HTML}{909090}
\definecolor{quarto-callout-note-color}{HTML}{0758E5}
\definecolor{quarto-callout-important-color}{HTML}{CC1914}
\definecolor{quarto-callout-warning-color}{HTML}{EB9113}
\definecolor{quarto-callout-tip-color}{HTML}{00A047}
\definecolor{quarto-callout-caution-color}{HTML}{FC5300}
\definecolor{quarto-callout-color-frame}{HTML}{acacac}
\definecolor{quarto-callout-note-color-frame}{HTML}{4582ec}
\definecolor{quarto-callout-important-color-frame}{HTML}{d9534f}
\definecolor{quarto-callout-warning-color-frame}{HTML}{f0ad4e}
\definecolor{quarto-callout-tip-color-frame}{HTML}{02b875}
\definecolor{quarto-callout-caution-color-frame}{HTML}{fd7e14}
\makeatother
\makeatletter
\@ifpackageloaded{caption}{}{\usepackage{caption}}
\AtBeginDocument{%
\ifdefined\contentsname
  \renewcommand*\contentsname{Table of contents}
\else
  \newcommand\contentsname{Table of contents}
\fi
\ifdefined\listfigurename
  \renewcommand*\listfigurename{List of Figures}
\else
  \newcommand\listfigurename{List of Figures}
\fi
\ifdefined\listtablename
  \renewcommand*\listtablename{List of Tables}
\else
  \newcommand\listtablename{List of Tables}
\fi
\ifdefined\figurename
  \renewcommand*\figurename{Figure}
\else
  \newcommand\figurename{Figure}
\fi
\ifdefined\tablename
  \renewcommand*\tablename{Table}
\else
  \newcommand\tablename{Table}
\fi
}
\@ifpackageloaded{float}{}{\usepackage{float}}
\floatstyle{ruled}
\@ifundefined{c@chapter}{\newfloat{codelisting}{h}{lop}}{\newfloat{codelisting}{h}{lop}[chapter]}
\floatname{codelisting}{Listing}
\newcommand*\listoflistings{\listof{codelisting}{List of Listings}}
\makeatother
\makeatletter
\makeatother
\makeatletter
\@ifpackageloaded{caption}{}{\usepackage{caption}}
\@ifpackageloaded{subcaption}{}{\usepackage{subcaption}}
\makeatother
\usepackage{bookmark}
\IfFileExists{xurl.sty}{\usepackage{xurl}}{} % add URL line breaks if available
\urlstyle{same}
\hypersetup{
  pdftitle={Scikit-longitudinal's TpT splitter: growing decision trees along the time dimension of longitudinal data},
  pdfauthor={Mathias Valla; Simon Provost},
  pdfkeywords={longitudinal data, decision
trees, reproducibility, Python, software},
  colorlinks=true,
  linkcolor={blue},
  filecolor={Maroon},
  citecolor={Blue},
  urlcolor={Blue},
  pdfcreator={LaTeX via pandoc}}


\title{Scikit-longitudinal's TpT splitter: growing decision trees along
the time dimension of longitudinal data}
\usepackage{etoolbox}
\makeatletter
\providecommand{\subtitle}[1]{% add subtitle to \maketitle
  \apptocmd{\@title}{\par {\large #1 \par}}{}{}
}
\makeatother
\subtitle{A short software note for Computo}
\author{Mathias Valla \and Simon Provost}
\date{2026-04-14}
\begin{document}
\maketitle
\begin{abstract}
We present the Time-penalized Trees (TpT) splitter as exposed in
scikit-longitudinal. TpT extends decision-tree split selection to
longitudinal settings by penalizing candidate splits that rely on later
observation waves. This note focuses on software usage and
reproducibility rather than a full methodological treatment. We document
the current TpT API in scikit-longitudinal and provide two compact
real-data examples: a subject-level classification task on the primary
biliary cirrhosis follow-up data and a subject-level regression task on
the dietox pig growth data. The accompanying repository contains a
frozen Python environment, data preparation scripts, and Quarto sources
for the paper.
\end{abstract}


\section{Introduction}\label{introduction}

\texttt{scikit-longitudinal} extends the \texttt{scikit-learn} ecosystem
with estimators and workflows tailored to repeated-measures data
(Provost and Freitas 2025). This software note focuses on the
Time-penalized Trees (TpT) splitter, a decision-tree split rule designed
for longitudinal covariates (Valla 2024). Rather than reproducing the
full methodological development, we document the current TpT API in the
package and provide two compact, fully reproducible real-data examples.
Beyond the original algorithmic presentation, TpT has also been shown to
be a consistent estimator for a non-empty set of observation times
(Valla and Milhaud 2026).

The repository is structured as required by Computo (Computo Team 2020):
the paper is written in Quarto, the environment is frozen in
\texttt{environment.yml}, and the data preparation and analysis code are
kept alongside the manuscript.

\section{\texorpdfstring{The TpT splitter in
\texttt{scikit-longitudinal}}{The TpT splitter in scikit-longitudinal}}\label{the-tpt-splitter-in-scikit-longitudinal}

At a tree node associated with parent time index \(t_p\), TpT scores a
candidate split at time \(t_c\) by multiplying the usual impurity
improvement \(\Delta I\) by an exponential time penalty:

\[
G_{\gamma} = \Delta I \exp \left[-\gamma (t_c - t_p)\right].
\]

The hyperparameter \(\gamma \ge 0\) controls how strongly later
measurements are penalized. When \(\gamma = 0\), TpT reduces to ordinary
impurity-based split selection. As \(\gamma\) increases, the tree
becomes more reluctant to use later waves unless they bring a clearly
stronger impurity reduction.

From a user perspective, the main entry points are
\texttt{TpTDecisionTreeClassifier} and
\texttt{TpTDecisionTreeRegressor}. Both can consume LONG-format data
directly by specifying an identifier column, an observation-time column
and a subject-specific duration column.

For this preliminary software note, we select \(\gamma\) on a small
candidate grid by minimizing the fitted tree's weighted total leaf
impurity on the training split. This is intentionally simple and should
be viewed as a placeholder for a more formal tuning strategy in the
final submission.

{Note to self: a Computo reviewer may reasonably ask for a stronger
justification of this gamma-selection rule, or for a validation-based
alternative. Keep the placeholder wording explicit and avoid overselling
it.}

\begin{tcolorbox}[enhanced jigsaw, titlerule=0mm, rightrule=.15mm, coltitle=black, opacitybacktitle=0.6, leftrule=.75mm, colback=white, bottomtitle=1mm, toptitle=1mm, opacityback=0, title=\textcolor{quarto-callout-note-color}{\faInfo}\hspace{0.5em}{Current scope}, arc=.35mm, breakable, bottomrule=.15mm, toprule=.15mm, left=2mm, colframe=quarto-callout-note-color-frame, colbacktitle=quarto-callout-note-color!10!white]

The present implementation is suitable for compact reproducible examples
and software-oriented demonstrations. Missing-value routing inside the
Cython TpT splitter is still treated as future work, and monotonic
constraints are explicitly disabled for TpT at the wrapper level.

\end{tcolorbox}

\section{Reproducible setup}\label{reproducible-setup}

The repository contains:

\begin{itemize}
\tightlist
\item
  \texttt{scripts/prepare\_example\_data.py} to download and preprocess
  the open datasets;
\item
  \texttt{analysis/tpt\_examples.py} with helper functions used by the
  paper;
\item
  \texttt{environment.yml} with a frozen Python environment.
\end{itemize}

The environment currently pins the TpT-enabled forks by GitHub SHA so
that the examples remain runnable before the upstream releases catch up
with the merged pull requests.

\section{\texorpdfstring{Example 1: Subject-level classification on
\texttt{pbcseq}}{Example 1: Subject-level classification on pbcseq}}\label{example-1-subject-level-classification-on-pbcseq}

Our classification example uses \texttt{pbcseq}, the sequential primary
biliary cirrhosis follow-up data distributed with the \texttt{survival}
package (Therneau, Terry 2025). We keep four complete longitudinal
biomarkers (\texttt{bili}, \texttt{albumin}, \texttt{protime},
\texttt{ast}), use \texttt{day} as the observation-time variable,
\texttt{futime} as the subject-specific horizon, and define a binary
target equal to one when the subject experiences either transplantation
or death by the end of follow-up.

The example is run on a subject-level train/test split, with the
long-format API of \texttt{TpTDecisionTreeClassifier}. The minimal
fit/predict workflow is:

\begin{Shaded}
\begin{Highlighting}[]
\ImportTok{from}\NormalTok{ scikit\_longitudinal.estimators.trees }\ImportTok{import}\NormalTok{ TpTDecisionTreeClassifier}

\NormalTok{feature\_columns }\OperatorTok{=}\NormalTok{ [}\StringTok{"id"}\NormalTok{, }\StringTok{"day"}\NormalTok{, }\StringTok{"futime"}\NormalTok{, }\StringTok{"bili"}\NormalTok{, }\StringTok{"albumin"}\NormalTok{, }\StringTok{"protime"}\NormalTok{, }\StringTok{"ast"}\NormalTok{]}

\NormalTok{clf }\OperatorTok{=}\NormalTok{ TpTDecisionTreeClassifier(}
\NormalTok{    gamma}\OperatorTok{=}\NormalTok{selected\_gamma,}
\NormalTok{    max\_depth}\OperatorTok{=}\DecValTok{3}\NormalTok{,}
\NormalTok{    min\_samples\_leaf}\OperatorTok{=}\DecValTok{10}\NormalTok{,}
\NormalTok{    random\_state}\OperatorTok{=}\DecValTok{42}\NormalTok{,}
\NormalTok{    assume\_long\_format}\OperatorTok{=}\VariableTok{True}\NormalTok{,}
\NormalTok{    id\_col}\OperatorTok{=}\StringTok{"id"}\NormalTok{,}
\NormalTok{    time\_col}\OperatorTok{=}\StringTok{"day"}\NormalTok{,}
\NormalTok{    duration\_col}\OperatorTok{=}\StringTok{"futime"}\NormalTok{,}
\NormalTok{    time\_step}\OperatorTok{=}\FloatTok{365.0}\NormalTok{,}
\NormalTok{    long\_feature\_columns}\OperatorTok{=}\NormalTok{[}\StringTok{"bili"}\NormalTok{, }\StringTok{"albumin"}\NormalTok{, }\StringTok{"protime"}\NormalTok{, }\StringTok{"ast"}\NormalTok{],}
\NormalTok{)}

\NormalTok{clf.fit(train\_df[feature\_columns], train\_df[}\StringTok{"event"}\NormalTok{])}
\NormalTok{pred }\OperatorTok{=}\NormalTok{ clf.predict(test\_df[feature\_columns])}
\NormalTok{proba }\OperatorTok{=}\NormalTok{ clf.predict\_proba(test\_df[feature\_columns])[:, }\DecValTok{1}\NormalTok{]}
\end{Highlighting}
\end{Shaded}

\begin{Shaded}
\begin{Highlighting}[]
\ImportTok{import}\NormalTok{ pandas }\ImportTok{as}\NormalTok{ pd}

\ImportTok{from}\NormalTok{ analysis.tpt\_examples }\ImportTok{import}\NormalTok{ run\_classification\_example}

\NormalTok{pbcseq\_example }\OperatorTok{=}\NormalTok{ run\_classification\_example()}
\NormalTok{pd.DataFrame([pbcseq\_example.metrics]).}\BuiltInTok{round}\NormalTok{(}\DecValTok{3}\NormalTok{)}
\end{Highlighting}
\end{Shaded}

\begin{verbatim}
/Users/mathiasvalla/Desktop/Projet_TPT/PR SCIKIT-LONG/.venv/lib/python3.10/site-packages/stopit/__init__.py:10: UserWarning: pkg_resources is deprecated as an API. See https://setuptools.pypa.io/en/latest/pkg_resources.html. The pkg_resources package is slated for removal as early as 2025-11-30. Refrain from using this package or pin to Setuptools<81.
  import pkg_resources
\end{verbatim}

\begin{longtable}[]{@{}lllllllllll@{}}

\caption{\label{tbl-pbcseq-metrics}Subject-level test metrics for the
pbcseq classification example.}

\tabularnewline

\toprule\noalign{}
& train\_subjects & test\_subjects & selected\_gamma & accuracy &
balanced\_accuracy & roc\_auc & tree\_total\_impurity &
tree\_mean\_impurity & tree\_depth & n\_nodes \\
\midrule\noalign{}
\endhead
\bottomrule\noalign{}
\endlastfoot
0 & 234 & 78 & 0.2 & 0.846 & 0.865 & 0.912 & 93.675 & 0.207 & 3 & 13 \\

\end{longtable}

We also report the small hyperparameter set used for this preliminary
example.

\begin{Shaded}
\begin{Highlighting}[]
\NormalTok{pbcseq\_example.hyperparameters}
\end{Highlighting}
\end{Shaded}

\begin{longtable}[]{@{}lll@{}}

\caption{\label{tbl-pbcseq-hyperparameters}Hyperparameter choices for
the pbcseq classification example.}

\tabularnewline

\toprule\noalign{}
& parameter & value \\
\midrule\noalign{}
\endhead
\bottomrule\noalign{}
\endlastfoot
0 & input\_format & long \\
1 & criterion & gini \\
2 & candidate\_gammas & 0.000 to 0.400 by 0.005 \\
3 & selected\_gamma & 0.2 \\
4 & max\_depth & 3 \\
5 & min\_samples\_leaf & 10 \\
6 & time\_step\_days & 365.0 \\
7 & selection\_rule & minimize training tree\_total\_impurity on a
0.0... \\

\end{longtable}

The selected gamma comes from the training-grid search shown in
Table~\ref{tbl-pbcseq-gamma-search}.

\begin{Shaded}
\begin{Highlighting}[]
\NormalTok{pbcseq\_example.gamma\_table.}\BuiltInTok{round}\NormalTok{(}\DecValTok{4}\NormalTok{)}
\end{Highlighting}
\end{Shaded}

\begin{longtable}[]{@{}llllll@{}}

\caption{\label{tbl-pbcseq-gamma-search}Gamma search for the pbcseq
classification example. The search uses a 0.005 grid, but the table
keeps only distinct training tree-total-impurity values and breaks ties
by retaining the smallest gamma.}

\tabularnewline

\toprule\noalign{}
& gamma & tree\_total\_impurity & tree\_mean\_impurity & tree\_depth &
n\_nodes \\
\midrule\noalign{}
\endhead
\bottomrule\noalign{}
\endlastfoot
0 & 0.000 & 204.5029 & 0.4388 & 3 & 9 \\
1 & 0.005 & 174.0299 & 0.4095 & 3 & 11 \\
2 & 0.010 & 221.3713 & 0.4730 & 3 & 11 \\
3 & 0.015 & 161.7476 & 0.3547 & 3 & 11 \\
4 & 0.060 & 194.9955 & 0.3645 & 3 & 11 \\
5 & 0.080 & 153.6755 & 0.2825 & 3 & 11 \\
6 & 0.115 & 94.5752 & 0.2189 & 3 & 11 \\
7 & 0.145 & 111.4667 & 0.2488 & 3 & 11 \\
8 & 0.190 & 228.8063 & 0.3945 & 3 & 13 \\
9 & 0.195 & 147.1358 & 0.2093 & 3 & 13 \\
10 & 0.200 & 93.6752 & 0.2068 & 3 & 13 \\
11 & 0.230 & 122.0084 & 0.2118 & 3 & 11 \\
12 & 0.255 & 119.3562 & 0.1871 & 3 & 11 \\
13 & 0.375 & 130.4875 & 0.1977 & 3 & 11 \\
14 & 0.385 & 132.9564 & 0.2277 & 3 & 13 \\
15 & 0.390 & 140.2477 & 0.2474 & 3 & 13 \\

\end{longtable}

The resulting subject-level metrics are shown in
Table~\ref{tbl-pbcseq-metrics}. A small preview of the held-out
predictions is given in Table~\ref{tbl-pbcseq-predictions}.

\begin{Shaded}
\begin{Highlighting}[]
\NormalTok{pbcseq\_example.predictions.head(}\DecValTok{8}\NormalTok{).}\BuiltInTok{round}\NormalTok{(}\DecValTok{3}\NormalTok{)}
\end{Highlighting}
\end{Shaded}

\begin{longtable}[]{@{}lllll@{}}

\caption{\label{tbl-pbcseq-predictions}Preview of held-out subject-level
predictions for the pbcseq classification example.}

\tabularnewline

\toprule\noalign{}
& id & observed & predicted & predicted\_probability \\
\midrule\noalign{}
\endhead
\bottomrule\noalign{}
\endlastfoot
22 & 78 & 1 & 1 & 1.000 \\
27 & 105 & 1 & 1 & 1.000 \\
70 & 278 & 1 & 1 & 1.000 \\
73 & 292 & 1 & 1 & 1.000 \\
74 & 293 & 0 & 1 & 1.000 \\
8 & 26 & 1 & 1 & 0.944 \\
26 & 94 & 1 & 1 & 0.944 \\
32 & 117 & 1 & 1 & 0.944 \\

\end{longtable}

The fitted tree is displayed below with the dedicated
\texttt{plot\_tree} exporter, using the wide feature names built
internally by the TpT preprocessing step.

\begin{Shaded}
\begin{Highlighting}[]
\ImportTok{from}\NormalTok{ sklearn }\ImportTok{import}\NormalTok{ tree}
\ImportTok{import}\NormalTok{ matplotlib.pyplot }\ImportTok{as}\NormalTok{ plt}

\NormalTok{fig, ax }\OperatorTok{=}\NormalTok{ plt.subplots(figsize}\OperatorTok{=}\NormalTok{(}\DecValTok{10}\NormalTok{, }\DecValTok{6}\NormalTok{))}
\NormalTok{tree.plot\_tree(}
\NormalTok{    pbcseq\_example.model,}
\NormalTok{    feature\_names}\OperatorTok{=}\NormalTok{pbcseq\_example.feature\_names,}
\NormalTok{    filled}\OperatorTok{=}\VariableTok{True}\NormalTok{,}
\NormalTok{    rounded}\OperatorTok{=}\VariableTok{True}\NormalTok{,}
\NormalTok{    fontsize}\OperatorTok{=}\DecValTok{9}\NormalTok{,}
\NormalTok{    tpt\_time\_scale}\OperatorTok{=}\FloatTok{0.5}\NormalTok{,}
\NormalTok{    ax}\OperatorTok{=}\NormalTok{ax,}
\NormalTok{)}
\NormalTok{ax.set\_title(}\SpecialStringTok{f"TpT classification tree (selected gamma = }\SpecialCharTok{\{}\NormalTok{pbcseq\_example}\SpecialCharTok{.}\NormalTok{metrics[}\StringTok{\textquotesingle{}selected\_gamma\textquotesingle{}}\NormalTok{]}\SpecialCharTok{\}}\SpecialStringTok{)"}\NormalTok{)}
\NormalTok{plt.show()}
\end{Highlighting}
\end{Shaded}

\begin{figure}[H]

\centering{

\pandocbounded{\includegraphics[keepaspectratio]{paper_files/figure-latex/fig-pbcseq-tree-output-1.pdf}}

}

\caption{\label{fig-pbcseq-tree}Selected TpT classification tree on
pbcseq, displayed with the dedicated plot\_tree exporter.}

\end{figure}%

This is precisely the kind of output we want in a software note: the
fitted tree shows explicitly which wave-specific biomarkers were
selected by TpT, while the calling interface remains close to ordinary
\texttt{scikit-learn} usage.

{Note to self: add one short sentence explaining why the selected
classification tree is clinically plausible, otherwise the figure risks
feeling purely decorative.}

\section{\texorpdfstring{Example 2: Subject-level regression on
\texttt{dietox}}{Example 2: Subject-level regression on dietox}}\label{example-2-subject-level-regression-on-dietox}

Our regression example uses \texttt{dietox}, a pig-growth longitudinal
dataset distributed with \texttt{geepack} (Halekoh, Ulrich and
Højsgaard, Søren and Yan, Jun 2025). For each pig, we hold out the last
recorded weight as the response \texttt{final\_weight}, and use the
earlier repeated weights as the longitudinal covariate history. This
produces a simple real-data regression task: predict the final observed
weight of a subject from its earlier growth trajectory.

In this example we explicitly use the TpT conversion helper
\texttt{long\_to\_wide(...)}, then fit the regressor on the resulting
wide matrix and its \texttt{feature\_groups} metadata. The minimal
workflow is:

\begin{Shaded}
\begin{Highlighting}[]
\ImportTok{from}\NormalTok{ scikit\_longitudinal.estimators.trees }\ImportTok{import}\NormalTok{ TpTDecisionTreeRegressor}
\ImportTok{from}\NormalTok{ scikit\_longitudinal.estimators.trees.TpT.\_preprocessing }\ImportTok{import}\NormalTok{ long\_to\_wide}

\NormalTok{feature\_columns }\OperatorTok{=}\NormalTok{ [}\StringTok{"Pig"}\NormalTok{, }\StringTok{"Time"}\NormalTok{, }\StringTok{"duration"}\NormalTok{, }\StringTok{"Weight"}\NormalTok{]}

\NormalTok{wide\_train }\OperatorTok{=}\NormalTok{ long\_to\_wide(}
\NormalTok{    train\_df[feature\_columns],}
\NormalTok{    id\_col}\OperatorTok{=}\StringTok{"Pig"}\NormalTok{,}
\NormalTok{    time\_col}\OperatorTok{=}\StringTok{"Time"}\NormalTok{,}
\NormalTok{    duration\_col}\OperatorTok{=}\StringTok{"duration"}\NormalTok{,}
\NormalTok{    time\_step}\OperatorTok{=}\FloatTok{1.0}\NormalTok{,}
\NormalTok{    feature\_columns}\OperatorTok{=}\NormalTok{[}\StringTok{"Weight"}\NormalTok{],}
\NormalTok{    target}\OperatorTok{=}\NormalTok{train\_df[}\StringTok{"final\_weight"}\NormalTok{],}
\NormalTok{)}

\NormalTok{reg }\OperatorTok{=}\NormalTok{ TpTDecisionTreeRegressor(}
\NormalTok{    gamma}\OperatorTok{=}\NormalTok{selected\_gamma,}
\NormalTok{    criterion}\OperatorTok{=}\StringTok{"squared\_error"}\NormalTok{,}
\NormalTok{    max\_depth}\OperatorTok{=}\DecValTok{3}\NormalTok{,}
\NormalTok{    min\_samples\_leaf}\OperatorTok{=}\DecValTok{5}\NormalTok{,}
\NormalTok{    random\_state}\OperatorTok{=}\DecValTok{42}\NormalTok{,}
\NormalTok{    features\_group}\OperatorTok{=}\NormalTok{wide\_train.feature\_groups,}
\NormalTok{)}

\NormalTok{reg.fit(wide\_train.X, wide\_train.y)}
\NormalTok{pred }\OperatorTok{=}\NormalTok{ reg.predict(wide\_test.X)}
\end{Highlighting}
\end{Shaded}

The next two tables show the same training subjects before and after the
explicit long-to-wide conversion.

\begin{Shaded}
\begin{Highlighting}[]
\ImportTok{from}\NormalTok{ analysis.tpt\_examples }\ImportTok{import}\NormalTok{ run\_regression\_example}

\NormalTok{dietox\_example }\OperatorTok{=}\NormalTok{ run\_regression\_example()}
\NormalTok{dietox\_example.long\_preview.}\BuiltInTok{round}\NormalTok{(}\DecValTok{2}\NormalTok{)}
\end{Highlighting}
\end{Shaded}

\begin{longtable}[]{@{}llllll@{}}

\caption{\label{tbl-dietox-long-preview}Preview of the long-format
dietox training data.}

\tabularnewline

\toprule\noalign{}
& Pig & Time & duration & Weight & final\_weight \\
\midrule\noalign{}
\endhead
\bottomrule\noalign{}
\endlastfoot
0 & 4602 & 1 & 11 & 28.3 & 106.7 \\
1 & 4602 & 2 & 11 & 30.1 & 106.7 \\
2 & 4602 & 3 & 11 & 38.3 & 106.7 \\
3 & 4602 & 4 & 11 & 44.5 & 106.7 \\
4 & 4602 & 5 & 11 & 51.6 & 106.7 \\
5 & 4602 & 6 & 11 & 57.6 & 106.7 \\
6 & 4602 & 7 & 11 & 65.0 & 106.7 \\
7 & 4602 & 8 & 11 & 73.0 & 106.7 \\

\end{longtable}

\begin{Shaded}
\begin{Highlighting}[]
\NormalTok{dietox\_example.wide\_preview.}\BuiltInTok{round}\NormalTok{(}\DecValTok{2}\NormalTok{)}
\end{Highlighting}
\end{Shaded}

\begin{longtable}[]{@{}lllllll@{}}

\caption{\label{tbl-dietox-wide-preview}Preview of the wide-format
matrix produced by long\_to\_wide for the dietox example.}

\tabularnewline

\toprule\noalign{}
& subject\_id & Weight\_t0 & Weight\_t1 & Weight\_t2 & Weight\_t3 &
final\_weight \\
\midrule\noalign{}
\endhead
\bottomrule\noalign{}
\endlastfoot
0 & 4602 & 28.3 & 28.3 & 30.1 & 38.3 & 106.7 \\
1 & 4603 & 27.6 & 27.6 & 30.6 & 38.7 & 105.4 \\
2 & 4605 & 31.5 & 31.5 & 34.8 & 40.7 & 109.2 \\
3 & 4645 & 27.7 & 27.7 & 33.6 & 44.0 & 102.3 \\
4 & 4756 & 17.0 & 17.0 & 19.0 & 23.6 & 89.3 \\
5 & 4757 & 26.9 & 26.9 & 32.3 & 38.6 & 91.0 \\
6 & 4813 & 27.3 & 27.3 & 32.4 & 38.3 & 99.8 \\
7 & 4815 & 24.9 & 24.9 & 29.7 & 33.8 & 102.7 \\

\end{longtable}

\begin{Shaded}
\begin{Highlighting}[]
\NormalTok{pd.DataFrame([dietox\_example.metrics]).}\BuiltInTok{round}\NormalTok{(}\DecValTok{3}\NormalTok{)}
\end{Highlighting}
\end{Shaded}

\begin{longtable}[]{@{}llllllllll@{}}

\caption{\label{tbl-dietox-metrics}Subject-level test metrics for the
dietox regression example.}

\tabularnewline

\toprule\noalign{}
& train\_subjects & test\_subjects & selected\_gamma & mae & r2 &
tree\_total\_impurity & tree\_mean\_impurity & tree\_depth & n\_nodes \\
\midrule\noalign{}
\endhead
\bottomrule\noalign{}
\endlastfoot
0 & 54 & 18 & 0.0 & 4.444 & 0.69 & 424.832 & 6.852 & 3 & 13 \\

\end{longtable}

\begin{Shaded}
\begin{Highlighting}[]
\NormalTok{dietox\_example.hyperparameters}
\end{Highlighting}
\end{Shaded}

\begin{longtable}[]{@{}lll@{}}

\caption{\label{tbl-dietox-hyperparameters}Hyperparameter choices for
the dietox regression example.}

\tabularnewline

\toprule\noalign{}
& parameter & value \\
\midrule\noalign{}
\endhead
\bottomrule\noalign{}
\endlastfoot
0 & input\_format & wide (after explicit long\_to\_wide conversion) \\
1 & criterion & squared\_error (MSE) \\
2 & candidate\_gammas & 0.000 to 0.200 by 0.005 \\
3 & selected\_gamma & 0.0 \\
4 & max\_depth & 3 \\
5 & min\_samples\_leaf & 5 \\
6 & time\_step & 1.0 \\
7 & selection\_rule & minimize training tree\_total\_impurity on a
0.0... \\

\end{longtable}

\begin{Shaded}
\begin{Highlighting}[]
\NormalTok{dietox\_example.gamma\_table.}\BuiltInTok{round}\NormalTok{(}\DecValTok{4}\NormalTok{)}
\end{Highlighting}
\end{Shaded}

\begin{longtable}[]{@{}llllll@{}}

\caption{\label{tbl-dietox-gamma-search}Gamma search for the dietox
regression example. The search uses a 0.005 grid, but the table keeps
only distinct training tree-total-impurity values and breaks ties by
retaining the smallest gamma.}

\tabularnewline

\toprule\noalign{}
& gamma & tree\_total\_impurity & tree\_mean\_impurity & tree\_depth &
n\_nodes \\
\midrule\noalign{}
\endhead
\bottomrule\noalign{}
\endlastfoot
0 & 0.000 & 424.8324 & 6.8521 & 3 & 13 \\
1 & 0.015 & 481.4917 & 7.7660 & 3 & 11 \\
2 & 0.075 & 453.6479 & 7.3169 & 3 & 11 \\
3 & 0.100 & 629.9755 & 10.1609 & 3 & 9 \\
4 & 0.120 & 708.6757 & 11.4303 & 3 & 9 \\

\end{longtable}

\begin{Shaded}
\begin{Highlighting}[]
\NormalTok{dietox\_example.predictions.head(}\DecValTok{8}\NormalTok{).}\BuiltInTok{round}\NormalTok{(}\DecValTok{2}\NormalTok{)}
\end{Highlighting}
\end{Shaded}

\begin{longtable}[]{@{}lllll@{}}

\caption{\label{tbl-dietox-predictions}Preview of held-out subject-level
predictions for the dietox regression example.}

\tabularnewline

\toprule\noalign{}
& Pig & observed & predicted & abs\_error \\
\midrule\noalign{}
\endhead
\bottomrule\noalign{}
\endlastfoot
4 & 4760 & 65.8 & 84.42 & 18.62 \\
6 & 4858 & 107.1 & 92.70 & 14.40 \\
10 & 5862 & 93.2 & 84.42 & 8.78 \\
1 & 4641 & 117.0 & 109.16 & 7.84 \\
13 & 6433 & 115.4 & 109.16 & 6.24 \\
14 & 6909 & 103.2 & 98.36 & 4.84 \\
17 & 8192 & 113.0 & 109.16 & 3.84 \\
8 & 5578 & 105.7 & 109.16 & 3.46 \\

\end{longtable}

\begin{Shaded}
\begin{Highlighting}[]
\ImportTok{import}\NormalTok{ matplotlib.pyplot }\ImportTok{as}\NormalTok{ plt}

\NormalTok{palette }\OperatorTok{=}\NormalTok{ \{}
    \StringTok{"Baseline history"}\NormalTok{: }\StringTok{"\#1b9e77"}\NormalTok{,}
    \StringTok{"Mid history"}\NormalTok{: }\StringTok{"\#d95f02"}\NormalTok{,}
    \StringTok{"Full history"}\NormalTok{: }\StringTok{"\#7570b3"}\NormalTok{,}
\NormalTok{\}}

\NormalTok{fig, ax }\OperatorTok{=}\NormalTok{ plt.subplots(figsize}\OperatorTok{=}\NormalTok{(}\FloatTok{6.5}\NormalTok{, }\FloatTok{4.5}\NormalTok{))}
\ControlFlowTok{for}\NormalTok{ stage, frame }\KeywordTok{in}\NormalTok{ dietox\_example.stage\_predictions.groupby(}\StringTok{"stage"}\NormalTok{, sort}\OperatorTok{=}\VariableTok{False}\NormalTok{):}
\NormalTok{    ax.scatter(}
\NormalTok{        frame[}\StringTok{"observed"}\NormalTok{],}
\NormalTok{        frame[}\StringTok{"predicted"}\NormalTok{],}
\NormalTok{        alpha}\OperatorTok{=}\FloatTok{0.75}\NormalTok{,}
\NormalTok{        label}\OperatorTok{=}\NormalTok{stage,}
\NormalTok{        color}\OperatorTok{=}\NormalTok{palette[stage],}
\NormalTok{    )}
\NormalTok{limits }\OperatorTok{=}\NormalTok{ [}
\NormalTok{    dietox\_example.stage\_predictions[[}\StringTok{"observed"}\NormalTok{, }\StringTok{"predicted"}\NormalTok{]].}\BuiltInTok{min}\NormalTok{().}\BuiltInTok{min}\NormalTok{(),}
\NormalTok{    dietox\_example.stage\_predictions[[}\StringTok{"observed"}\NormalTok{, }\StringTok{"predicted"}\NormalTok{]].}\BuiltInTok{max}\NormalTok{().}\BuiltInTok{max}\NormalTok{(),}
\NormalTok{]}
\NormalTok{ax.plot(limits, limits, linestyle}\OperatorTok{=}\StringTok{"{-}{-}"}\NormalTok{, linewidth}\OperatorTok{=}\DecValTok{1}\NormalTok{, color}\OperatorTok{=}\StringTok{"black"}\NormalTok{)}
\NormalTok{ax.set\_xlabel(}\StringTok{"Observed final weight"}\NormalTok{)}
\NormalTok{ax.set\_ylabel(}\StringTok{"Predicted final weight"}\NormalTok{)}
\NormalTok{ax.set\_title(}\StringTok{"TpT regression on dietox"}\NormalTok{)}
\NormalTok{ax.legend(frameon}\OperatorTok{=}\VariableTok{False}\NormalTok{, loc}\OperatorTok{=}\StringTok{"upper left"}\NormalTok{)}
\NormalTok{plt.show()}
\end{Highlighting}
\end{Shaded}

\begin{figure}[H]

\centering{

\pandocbounded{\includegraphics[keepaspectratio]{paper_files/figure-latex/fig-dietox-predictions-output-1.pdf}}

}

\caption{\label{fig-dietox-predictions}Observed versus predicted final
pig weights in the dietox regression example, with predictions computed
from baseline-only, mid-history, and full-history inputs.}

\end{figure}%

The selected regression tree can also be shown directly with the
dedicated \texttt{plot\_tree} exporter.

\begin{Shaded}
\begin{Highlighting}[]
\NormalTok{fig, ax }\OperatorTok{=}\NormalTok{ plt.subplots(figsize}\OperatorTok{=}\NormalTok{(}\DecValTok{11}\NormalTok{, }\DecValTok{6}\NormalTok{))}
\NormalTok{tree.plot\_tree(}
\NormalTok{    dietox\_example.model,}
\NormalTok{    feature\_names}\OperatorTok{=}\NormalTok{dietox\_example.feature\_names,}
\NormalTok{    filled}\OperatorTok{=}\VariableTok{True}\NormalTok{,}
\NormalTok{    rounded}\OperatorTok{=}\VariableTok{True}\NormalTok{,}
\NormalTok{    fontsize}\OperatorTok{=}\DecValTok{9}\NormalTok{,}
\NormalTok{    tpt\_time\_scale}\OperatorTok{=}\FloatTok{0.30}\NormalTok{,}
\NormalTok{    ax}\OperatorTok{=}\NormalTok{ax,}
\NormalTok{)}
\NormalTok{ax.set\_title(}\SpecialStringTok{f"TpT regression tree (selected gamma = }\SpecialCharTok{\{}\NormalTok{dietox\_example}\SpecialCharTok{.}\NormalTok{metrics[}\StringTok{\textquotesingle{}selected\_gamma\textquotesingle{}}\NormalTok{]}\SpecialCharTok{\}}\SpecialStringTok{)"}\NormalTok{)}
\NormalTok{plt.show()}
\end{Highlighting}
\end{Shaded}

\begin{figure}[H]

\centering{

\pandocbounded{\includegraphics[keepaspectratio]{paper_files/figure-latex/fig-dietox-tree-output-1.pdf}}

}

\caption{\label{fig-dietox-tree}Selected TpT regression tree on dietox,
displayed with the dedicated plot\_tree exporter.}

\end{figure}%

{Note to self: the regression example is reproducible and useful for
software documentation, but a reviewer may ask whether a single-feature
longitudinal history is too simple for the final note. Keep a richer
backup dataset in mind.}

{Note to self: the visual style is still functional rather than
polished. Before submission, improve the figure styling and overall
visual consistency across the tree plots, scatter plot, and tables.}

\section{Discussion}\label{discussion}

These two examples are intentionally small. Their role is not to provide
a full benchmarking study, but to document a reproducible and
inspectable usage path for TpT within \texttt{scikit-longitudinal}.
Together they illustrate three software-facing points:

\begin{enumerate}
\def\labelenumi{\arabic{enumi}.}
\tightlist
\item
  subject-level train/test evaluation can be carried out while keeping
  the inputs in LONG format, or by explicit conversion to the wide
  matrix expected by the core TpT splitter;
\item
  TpT exposes a familiar \texttt{scikit-learn}-like interface through
  \texttt{fit}, \texttt{predict}, and \texttt{predict\_proba};
\item
  preliminary gamma tuning can be implemented with a compact
  impurity-based helper before moving to a more formal model-selection
  strategy;
\item
  the learned trees remain readable because the selected splits are tied
  to explicit wave-indexed feature names.
\end{enumerate}

Future revisions of this note can stay short while becoming more robust
by switching the environment from pinned GitHub SHAs to released package
versions, and by updating the examples once the remaining TpT
implementation limitations around missing-value routing are resolved.

{Note to self: explicitly state near submission time whether the paper
is tied to released package versions or to pinned development SHAs,
because Computo reviewers will care about long-term reproducibility.}

\section{Acknowledgements}\label{acknowledgements}

Work(s) conducted within the Research Chair DIALog under the aegis of
the Risk Foundation, a joint initiative by Université Paris-Dauphine PSL
and CNP Assurances.

Placeholder for Simon Provost's acknowledgements.

{Note to self: fill Simon's affiliation and acknowledgements before
submission, and check whether Computo wants a software-availability or
funding sentence in a more specific location.}

\section*{References}\label{references}
\addcontentsline{toc}{section}{References}

{Note to self: complete and clean the bibliography before submission,
including any missing software, package, and methodological references
cited in the final text.}

\phantomsection\label{refs}
\begin{CSLReferences}{1}{0}
\bibitem[\citeproctext]{ref-computo}
Computo Team. 2020. {``Computo: Reproducible Computational/Algorithmic
Contributions in Statistics and Machine Learning.''}

\bibitem[\citeproctext]{ref-dataset-dietox}
Halekoh, Ulrich and Højsgaard, Søren and Yan, Jun. 2025. {``Dietox:
Growth Curves of Pigs in a Factorial Experiment.''}
\url{https://vincentarelbundock.github.io/Rdatasets/doc/geepack/dietox.html}.

\bibitem[\citeproctext]{ref-provost2025sklong}
Provost, Simon, and Alex Freitas. 2025. {``Scikit-Longitudinal: A
Machine Learning Library for Longitudinal Classification in Python.''}
\emph{Journal of Open Source Software} 10 (112): 8481.
\url{https://doi.org/10.21105/joss.08481}.

\bibitem[\citeproctext]{ref-dataset-pbcseq}
Therneau, Terry. 2025. {``Pbcseq: Mayo Clinic Primary Biliary Cirrhosis
Sequential Data.''} \url{https://rdrr.io/cran/survival/man/pbcseq.html}.

\bibitem[\citeproctext]{ref-valla2024tpt}
Valla, Mathias. 2024. {``Time-Penalized Trees ({TpT}): Introducing a New
Tree-Based Data Mining Algorithm for Time-Varying Covariates.''}
\emph{Annals of Mathematics and Artificial Intelligence} 92: 1609--61.
\url{https://doi.org/10.1007/s10472-024-09950-w}.

\bibitem[\citeproctext]{ref-valla2026consistency}
Valla, Mathias, and Xavier Milhaud. 2026. {``Consistent Time-Aware Trees
for Longitudinal Data: The Time-Penalized Tree.''}
\url{https://cnrs.hal.science/hal-05022929}.

\end{CSLReferences}




\end{document}
